\documentclass[10pt]{article}
\providecommand{\FIGDIR}{fig}
\newcommand{\DIGESTDATE}{12 September 2026}
\newcommand{\PERIODSTART}{6 September 2026}
\newcommand{\PERIODEND}{12 September 2026}
% ===================== PM Research Watch - shared preamble =====================
\usepackage[T1]{fontenc}
\usepackage[utf8]{inputenc}
\usepackage{amsmath,amssymb}
\usepackage{tgpagella}
\usepackage{tgheros}
\renewcommand{\sfdefault}{qhv}
\renewcommand{\ttdefault}{lmtt}
\usepackage{microtype}
\usepackage[english]{babel}

\usepackage[a4paper,top=17mm,bottom=17mm,left=17mm,right=17mm,
            headsep=6mm,footskip=10mm]{geometry}
\usepackage{graphicx}
\usepackage{xcolor}
\usepackage{colortbl}
\usepackage{tikz}
\usetikzlibrary{calc}
\usepackage[most]{tcolorbox}
\usepackage{titlesec}
\usepackage{fancyhdr}
\usepackage{booktabs}
\usepackage{array}
\usepackage{longtable}
\usepackage{enumitem}
\usepackage{multicol}
\usepackage{ragged2e}
\usepackage[hidelinks]{hyperref}

% ------------------------------------------------------------------ palette
\definecolor{Ink}    {HTML}{14262E}
\definecolor{Deep}   {HTML}{0F4C5C}
\definecolor{Teal}   {HTML}{1F7A8C}
\definecolor{Sky}    {HTML}{4FA3C4}
\definecolor{Sage}   {HTML}{6FA287}
\definecolor{Amber}  {HTML}{E3A008}
\definecolor{Clay}   {HTML}{B5643C}
\definecolor{Coral}  {HTML}{D1495B}
\definecolor{Violet} {HTML}{6A5B9E}
\definecolor{Slate}  {HTML}{4A6572}
\definecolor{Paper}  {HTML}{FBF9F5}
\definecolor{Mist}   {HTML}{EEEAE2}
\definecolor{Rule}   {HTML}{D9D3C9}

\pagecolor{Paper}
\color{Ink}

\hypersetup{colorlinks=true, linkcolor=Deep, urlcolor=Teal, citecolor=Teal,
            pdftitle={PM Research Watch}, pdfauthor={Automated daily digest}}

% ------------------------------------------------------------------ headings
\titleformat{\section}
  {\sffamily\bfseries\large\color{Deep}}{}{0pt}
  {}[{\vspace{1.2mm}\color{Amber}\titlerule[1.1pt]}]
\titlespacing*{\section}{0pt}{5.5mm}{2.8mm}

\titleformat{\subsection}
  {\sffamily\bfseries\normalsize\color{Teal}}{}{0pt}{}
\titlespacing*{\subsection}{0pt}{3.6mm}{1.4mm}

\setlist[itemize]{leftmargin=4.2mm, itemsep=1.1mm, topsep=1.3mm, parsep=0pt,
                  label=\textcolor{Teal}{\footnotesize$\blacktriangleright$}}

% ------------------------------------------------------------------ header/footer
\pagestyle{fancy}
\fancyhf{}
\renewcommand{\headrulewidth}{0pt}
\renewcommand{\footrulewidth}{0pt}
\fancyhead[L]{\sffamily\scriptsize\color{Slate}PM RESEARCH WATCH\;\textcolor{Rule}{\textbar}\;\DIGESTDATE}
\fancyhead[R]{\sffamily\scriptsize\color{Slate}Particulate matter monitoring \& health effects}
\fancyfoot[C]{\sffamily\scriptsize\color{Slate}\thepage}

% ------------------------------------------------------------------ components
% metric tile
\newcommand{\metric}[3]{%
  \begin{tcolorbox}[enhanced, colback=white, colframe=Rule, boxrule=0.5pt,
    arc=1.6mm, left=2.2mm, right=2.2mm, top=1.8mm, bottom=1.8mm,
    borderline west={1.7mm}{0pt}{#3}, nobeforeafter, valign=center,
    height=17mm, width=\linewidth]
    \RaggedRight
    {\sffamily\bfseries\fontsize{17}{17}\selectfont\color{#3}#1}\\[0.5mm]
    {\sffamily\fontsize{6.5}{7.4}\selectfont\color{Slate}#2}
  \end{tcolorbox}}

% pull-quote / key insight box
\newtcolorbox{keybox}[1][]{enhanced, colback=Mist, colframe=Mist,
  arc=1.6mm, left=3.2mm, right=3.2mm, top=2.4mm, bottom=2.4mm,
  borderline west={1.7mm}{0pt}{Amber}, #1}

% subtopic band
\newcommand{\band}[2]{%
  \vspace{3.2mm}\noindent
  \begin{tcolorbox}[enhanced, colback=#2, colframe=#2, arc=1.2mm,
    left=2.6mm, right=2.6mm, top=1.1mm, bottom=1.1mm, boxrule=0pt,
    width=\linewidth, nobeforeafter]
    {\sffamily\bfseries\small\color{white}#1}
  \end{tcolorbox}\vspace{1.6mm}}

% one paper entry: {short cite}{title}{journal}{doi}{body}
\newcommand{\paperentry}[5]{%
  \par\vspace{1.9mm}\noindent
  {\sffamily\bfseries\footnotesize\color{Ink}#1}\;%
  {\sffamily\fontsize{7}{8}\selectfont\color{Slate}#3}\\[0.4mm]
  {\footnotesize\itshape\color{Deep}#2}\\[0.6mm]
  {\footnotesize #4}\\[0.3mm]
  {\fontsize{6.6}{7.6}\selectfont\color{Slate}\texttt{doi:}\,\href{https://doi.org/#5}{\texttt{\detokenize{#5}}}}
  \par}

% caption style
\newcommand{\figcap}[1]{\par\vspace{1.0mm}
  {\fontsize{7.2}{8.6}\selectfont\color{Slate}\textsf{\textbf{Fig.}}\; #1\par}\vspace{1.5mm}}

\providecommand{\permil}{\textperthousand}

\setlength{\parindent}{0pt}
\setlength{\parskip}{0pt}
\renewcommand{\arraystretch}{1.18}


\begin{document}

% ==================================================================== MASTHEAD
\thispagestyle{empty}
\begin{tikzpicture}[remember picture, overlay]
  \fill[Deep] (current page.north west) rectangle
      ($(current page.north east)+(0,-67mm)$);
  \fill[Teal] ($(current page.north west)+(0,-67mm)$) rectangle
      ($(current page.north east)+(0,-68.4mm)$);
  \foreach \x/\y/\r in {22/12/0.9, 41/26/1.5, 63/9/0.7, 88/31/1.1, 112/17/1.9,
                        138/29/0.8, 159/13/1.3, 178/54/1.0, 196/21/1.6,
                        30/44/0.6, 74/57/1.2, 125/48/0.9, 168/6/0.8, 191/59/0.7,
                        52/18/1.0, 100/60/0.6, 148/20/1.1, 66/49/1.4, 205/45/0.8}
    \fill[Sky, opacity=0.30] ($(current page.north west)+(\x mm,-\y mm)$) circle (\r mm);
\end{tikzpicture}

\vspace*{4mm}
{\sffamily\fontsize{7.6}{9}\selectfont\color{Sky}%
 WEEKLY ROLLUP\;\textbullet\;PhD RESEARCH WATCH\par}
\vspace{2.2mm}
{\sffamily\bfseries\fontsize{27}{29}\selectfont\color{white}%
 Particulate Matter\par}
{\sffamily\fontsize{27}{29}\selectfont\color{Sky}%
 Week in Review\par}
\vspace{5.0mm}
{\sffamily\fontsize{8.4}{10}\selectfont\color{white}%
 \mbox{Week ending \DIGESTDATE}\;\;\textcolor{Teal}{\textbar}\;\;\mbox{\PERIODSTART{} to \PERIODEND}\;\;
 \textcolor{Teal}{\textbar}\;\;\mbox{126 records pooled from 7 daily issues}\par}
\vspace{18mm}

% -------------------------------------------------------------- metric strip
\noindent
\begin{minipage}[t]{0.190\linewidth}\metric{126}{RECORDS POOLED\\(126 UNIQUE DOIs)}{Deep}\end{minipage}\hfill
\begin{minipage}[t]{0.190\linewidth}\metric{40\%}{MEASUREMENT SIDE\\(51 RECORDS)}{Teal}\end{minipage}\hfill
\begin{minipage}[t]{0.190\linewidth}\metric{10}{SUBTOPICS\\ACTIVE}{Sky}\end{minipage}\hfill
\begin{minipage}[t]{0.190\linewidth}\metric{47}{POOLED EFFECT\\ESTIMATES}{Amber}\end{minipage}\hfill
\begin{minipage}[t]{0.190\linewidth}\metric{20}{ABSTRACT-FREE\\(16\% OF WEEK)}{Clay}\end{minipage}

\vspace{6mm}

% -------------------------------------------------------------- week in one box
\section*{\sffamily\bfseries\color{Deep}The week in one paragraph}
\vspace{-1.5mm}{\color{Teal}\rule{\linewidth}{1.1pt}}\vspace{2.5mm}

\begin{keybox}
{\sffamily\bfseries\small\color{Deep}Sensing led a week for the first time since
mid-August, the quantitative epidemiology came back hard after last week's floor --- and
one in six records this week has no abstract in existence anywhere. The pipeline's
best-performing leg and its worst failure mode are now the same leg.}\\[1.4mm]
\footnotesize
\textbf{Scale.} 126 records over seven issues on daily totals of 15, 20, 18, 26, 16, 10
and \textbf{21}, a far flatter week than the 5.4$\times$ range of the last one. Against
recent rollups: 126 after 150, 109, 112 and 134 --- squarely mid-range, and the smoothness
matters more than the level, because week-over-week cluster movements below are not being
read through a single 38-record deposit day.

\textbf{\textit{Sensing, forecasting \& instrumentation} at 31 is the largest cluster, and
that is the week's structural news.} For three consecutive weeks \textit{Exposure
assessment \& modelling} had led; it has now halved, \textbf{47 $\rightarrow$ 20}, the
sharpest single-cluster contraction this watch has recorded, while Sensing fell only
34 $\rightarrow$ 31. The measurement side as a whole went \textbf{54\% $\rightarrow$
40\%}, so the two clusters did not move together: modelling-side exposure work
collapsed and instrumentation held.

\textbf{Estimate density recovered to one per 2.7 records.} 47 pooled estimates against
126 records, from a four-week series of 2.7, 2.9, 4.4, \textbf{2.7}. Last week's rollup
flagged 4.4 as the weakest on record and declined to say whether it was a real thinning
or a measurement-heavy artefact. One week later it is back to the series norm, which
settles it: \textbf{it was the artefact}. The week's quantitative spine is cardiometabolic
--- CHARLS coexposure, the Zhongshan AMI case-crossover, the Shijiazhuang AECOPD series,
the Hangzhou aortic-dissection component analysis, the Korean 8.25-million asthma cohort
and the Kadoorie fracture cohort.

\textbf{The abstract-free problem is now large enough to be a finding about the corpus.}
\textbf{20 of 126 records (16\%)} shipped as \textsc{metadata only}, up from 11\% last
week, and \textbf{10 of those 20 are in \textit{Sensing}} --- the pooled heatmap's densest
cell is \textit{Sensing} $\times$ \textit{Metadata only}. Crossref-by-ISSN is the only leg
that reaches \textit{J. Aerosol Sci.}, \textit{Atmos. Res.}, \textit{AMT} and \textit{ACP},
and it indexes a DOI days before the publisher deposits the abstract. So the leg that makes
the instrumentation literature visible at all is the same leg that delivers it unreadable.
\textbf{An abstract-retry pass over the standing list is the highest-value pipeline work
outstanding} --- it would convert roughly a sixth of the week from title-only to summarised.
\end{keybox}

\clearpage

% -------------------------------------------------------------- trend
\section*{\sffamily\bfseries\color{Deep}Subtopic trend and inflection}
\vspace{-1.5mm}{\color{Teal}\rule{\linewidth}{1.1pt}}\vspace{2.5mm}

\noindent\includegraphics[width=\linewidth]{\FIGDIR/w1_subtopic_trend.png}
\figcap{Per-issue subtopic counts across the week from \texttt{state/metrics.csv}, seven
issues. Daily totals were 15, 20, 18, 26, 16, 10 and 21 --- a 2.6$\times$ range against
5.4$\times$ last week, so the per-cluster lines are readable as movement rather than as
denominator noise. Two of the seven days are weekends (6 and 12 September) and both are
PubMed-suppressed: 12 September admitted \textbf{26 records registry-wide} against 1{,}135
on 8 September. Those two days were carried almost entirely by Europe PMC and
Crossref-by-ISSN, which is why the \textit{Sensing} line does not dip at the weekend the
way the health clusters do.}

\footnotesize
\textbf{Expanding.} \textit{Burden, policy \& mitigation} \textbf{7 $\rightarrow$ 14}
doubles off last week's recorded floor and is the week's largest proportional gain.
\textit{Respiratory \& allergic} \textbf{5 $\rightarrow$ 12} likewise more than doubles
from a floor. \textit{Neuro / mental health} \textbf{4 $\rightarrow$ 7} and
\textit{Cardiovascular \& metabolic} \textbf{7 $\rightarrow$ 8} both continue upward.
Read together, the health side recovered as a bloc from a single unusually thin week ---
which is the strongest evidence available that last week's health-side floor was a
sampling artefact and not a turn in the field.

\vspace{1.2mm}
\textbf{Contracting.} \textit{Exposure assessment \& modelling} \textbf{47 $\rightarrow$
20} is a halving from a cluster record, and no interpretation of it should outrun that
fact: a record high is the most likely thing to be followed by a fall. \textit{Occupational
\& indoor} \textbf{19 $\rightarrow$ 15} eases off last week's surge but stays well above
the pre-surge regular-week maximum of 9 --- \textbf{the indoor turn the 5 September rollup
declined to call a trend has now held for a second week}, which was the stated test.
\textit{Mechanistic toxicology} \textbf{11 $\rightarrow$ 9} and \textit{Sensing}
\textbf{34 $\rightarrow$ 31} are both mild.

\vspace{1.2mm}
\textbf{The inflection worth naming.} The four-week measurement-side share now reads
41, 47, 54, \textbf{40\%} --- it oscillates, it does not trend, and the 29 August rollup's
warning against building a claim on it stands. What has changed is \emph{within} the
measurement side: \textit{Exposure} and \textit{Sensing} moved in opposite directions
relative to each other for the first time in four weeks. One week is not a pattern. The
concrete thing to watch next week is whether \textit{Exposure} returns to the mid-20s,
which would confirm this week as the correction of a record rather than a shift.

%% LAYOUT RULE (2026-08-02): architecture donut in the whitespace here, not on its own page.
\vspace{2mm}
\noindent\includegraphics[width=\linewidth]{\FIGDIR/f2a_architecture.png}
\figcap{Study architecture over the pooled week. \textit{Measurement campaign} (23)
retakes the lead from \textit{Modelling / inventory} (19), reversing the first-ever
modelling-ahead-of-measurement week recorded in the last rollup. \textit{Metadata only} at
\textbf{20} is the second-largest architecture in the week, which is not an architecture at
all but an absence of one. Primary observational epidemiology (cohort 15,
cross-sectional 10, acute 6) totals 31, ahead of modelling for the first time in three
weeks. \textbf{Two} interventional records appear in the 126 --- a Dutch school
air-cleaner RCT and a smoke-free-home trial in public housing --- against one last week.}

\clearpage

% -------------------------------------------------------------- pooled analytics
\section*{\sffamily\bfseries\color{Deep}Pooled analytics}
\vspace{-1.5mm}{\color{Teal}\rule{\linewidth}{1.1pt}}\vspace{3mm}

\noindent\includegraphics[width=\linewidth]{\FIGDIR/f1_subtopics.png}
\figcap{Pooled subtopic distribution, 126 records. All ten clusters were active, for the
seventh consecutive week. The two measurement clusters (31 + 20 = 51) no longer outweigh
the eight health and policy clusters combined (75) --- last week they did, 81 to 69.}

\vspace{2mm}
\noindent\includegraphics[width=\linewidth]{\FIGDIR/f2b_endpoint.png}
\figcap{Health endpoint by organ system. \textbf{55 of 126 records carry no human
endpoint} --- monitoring (27), instrument (9), emissions (6), aerosol chemistry (5) and two
explicit none-labels --- a 44\% share, close to last week's 38\%. Among endpoints that
exist, \textit{Respiratory} at \textbf{21} is far ahead of anything else and is more than
the next three combined (mortality 5, cardiovascular 5, other clinical 5). Note that
\textit{Cancer} and \textit{Oncologic} remain split in \texttt{ENDPOINT\_CANON}, so
oncological work is understated by this panel; logged 3 September, still unfixed.}

\vspace{2mm}
\noindent\includegraphics[width=\linewidth]{\FIGDIR/f3_metric_geography.png}
\figcap{Left --- particle metric. Unspeciated PM and emission factors (52) lead
PM$_{2.5}$-only (39), with PM$_{2.5}$+PM$_{10}$ jointly at 18. \textbf{Ultrafine and
size-resolved number together total 3 of 126}, against seven records whose metric is an
optical property --- the field is measuring what instruments report, not what the source
signal lives in. Right --- geographic provenance. \textbf{China supplies 32 of 126 (25\%)},
with the \textit{Global / multi-region} fallback at 31, Europe 17, North America 14 and
East Asia (ex-China) 10. Sub-Saharan Africa at 7 is the strongest week this archive has
recorded for the region. The fallback bar is inflated by laboratory and computational
records that have no site by construction, and by daily-issue geography labels that were
never normalised --- that defect is carried and is logged again below.}

\clearpage

\noindent\includegraphics[width=\linewidth]{\FIGDIR/f5_forest.png}
\vspace{1mm}
\noindent\includegraphics[width=\linewidth]{\FIGDIR/f5_forest_2.png}
\figcap{All 47 pooled effect estimates on a common logarithmic axis, split across two
panels. \textbf{These are not mutually comparable and are not meta-analysed}: contrasts
include per-IQR, per-10\,$\mu$g\,m$^{-3}$, per-quartile, per-SD, per-smoke-day and
mixture-index units, on OR, HR, RR, IRR, PR and RRR. \textbf{Heterogeneity is structural,
not statistical}, so no $I^2$ is reported; pooling these would be a category error. Three
things are readable anyway. First, \textbf{four estimates strictly cross the null} --- the
CLHLS SO$_2$ mortality hazard, the Alzheimer's and vascular-dementia sex-modification
ratios, and the smoke-free-home 90-day primary outcome --- with two more whose lower bound
sits exactly at 1.00 (ozone--myopia, NO$_2$--multimorbidity), against two of 34 last week.
A further \textbf{seven estimates carry no interval at all} because the abstract reported
none; they are drawn as bare points and should not be read as precise. Second, \textbf{the
large point estimates are the imprecise ones}: 3.83 (2.63--5.57) for smoke-free-home
adoption and 2.51 (1.62--3.88) for acute stress, while the tightest interval on the axis
(1.33, 1.32--1.34, Korean incident asthma) belongs to one of the smallest effects and comes
from 8.25 million people. Third, several
estimates here are \emph{co-pollutant} rather than PM --- ozone, NO$_2$, SO$_2$ --- and are
plotted because the papers report them alongside PM; read the labels before reading the
positions.}

\vspace{2mm}
\noindent\includegraphics[width=\linewidth]{\FIGDIR/f6_lifecourse.png}
\figcap{Life-course windows across the pooled week. Working-age adulthood dominates as
usual; the change against last week is childhood, lifted by two measurement-side records
(schoolchild personal monitoring, a Dutch school air-cleaner trial) rather than by
epidemiology.}

\vspace{2mm}
\noindent\includegraphics[width=\linewidth]{\FIGDIR/f4_heatmap.png}
\figcap{Subtopic $\times$ architecture, pooled into 51 cells. The densest cell in the
entire week is \textit{Sensing} $\times$ \textit{Metadata only} at \textbf{10} --- ahead of
\textit{Sensing} $\times$ \textit{Measurement campaign} (9) and \textit{Sensing} $\times$
\textit{Modelling / inventory} (9). The standing structural finding survives a seventh
week: of the \textbf{51 records} in the \textit{Sensing} and \textit{Exposure} rows,
\textbf{5 carry any human health endpoint}. Instrumentation and epidemiology are still
being written by different people in different journals, and this week the instrumentation
half arrived without abstracts.}

\clearpage

% -------------------------------------------------------------- venues + trials
\section*{\sffamily\bfseries\color{Deep}Venues, trial watch, and what the pipeline did}
\vspace{-1.5mm}{\color{Teal}\rule{\linewidth}{1.1pt}}\vspace{2.5mm}

\footnotesize
\textbf{Venues.} Far less concentrated than last week: \textit{Environ Pollut}
\textbf{10 of 126} leads, against \textit{Environ Sci Technol}'s 23 of 150 a week ago.
\textit{Front Public Health} (7), \textit{Environ Res} (6) and \textit{Atmos Chem Phys}
(6) follow, then \textit{Atmosphere} (5) and \textit{J Aerosol Sci} (4). The venue table
is a harvesting artefact as much as an editorial one --- titles on the Crossref by-ISSN
list enter at their own publication rate, while journals reached only through PubMed enter
at PubMed's cadence, and this week PubMed's cadence included two near-dead weekend days.
\textbf{Preprints supplied 8 records} --- 4 Research Square, 4 Preprints.org --- all
labelled as such in the register.

\vspace{1.5mm}
\textbf{Recurring authors.} Unchanged and unresolvable: surname clustering is dominated by
high-frequency Chinese surnames, which is a naming artefact rather than a research-group
signal, and this watch still has no author-disambiguation or affiliation field.
\textbf{No repeat-group cluster is claimed for this week.} Adding a first-author
affiliation field remains the open item that would make this section mean anything; it has
been on the list since 3 September and is now the oldest unaddressed schema gap.

\vspace{1.5mm}
\textbf{Trial watch --- a genuine zero week, the first this watch has recorded.}
ClinicalTrials.gov returned \textbf{0 PM or air-pollution registrations or updates across
6--12 September}, against \textbf{4{,}413 studies updated registry-wide} in the same
window. That control matters because the 5 September rollup established that weekend zeros
are \emph{structural} --- the registry posts nothing on Saturdays and Sundays --- and could
not, on weekend data alone, distinguish a quiet registry from a quiet field. A seven-day
window containing five active weekdays and 4{,}413 registry-wide updates removes that
ambiguity: \textbf{this is a topical zero}. Against one new interventional registration
and two updates last week, it is a real if unremarkable fluctuation, and one week of it
means nothing on its own. \texttt{trials.json} still holds \textbf{26 trials}, now across
\textbf{12 windows}. \texttt{analyze\_endpoints} was not run --- there was nothing to run
it on.

\vspace{1.5mm}
\textbf{What the harvesting legs did.} \textbf{Crossref-by-ISSN was the week's workhorse
and its liability}, carrying the instrumentation literature PubMed cannot see and
delivering 10 of its \textit{Sensing} records without abstracts. \textbf{PubMed} carried
five of seven days and was effectively absent on both weekend days --- 26 records
registry-wide on 12 September is the lowest single-day figure this watch has measured, and
the 6 September Sunday behaved the same way. \textbf{Europe PMC} was the quiet success,
returning 17 on 12 September alone including two Research Square preprints, and it is now
the leg that keeps weekend issues from collapsing. \textbf{OpenAlex returned HTTP 429 on
every day of the week}, unchanged since 28 July, and remains a dead leg. \textbf{arXiv}
returned 429 or nothing in window all week and contributed no record to any issue in this
period --- second consecutive week of that, and the query is now overdue for inspection.
\textbf{Consensus returned a nil every day}, and from 9 September onward the nil has been
verified the correct way, on Crossref \texttt{created} date rather than publication date;
its hits are consistently the right literature with dates years outside the window.

\vspace{2mm}
% -------------------------------------------------------------- gaps
\section*{\sffamily\bfseries\color{Deep}Open gaps}
\vspace{-1.5mm}{\color{Teal}\rule{\linewidth}{1.1pt}}\vspace{2.5mm}

\footnotesize
\begin{itemize}
\item \textbf{Twenty records this week assert no finding.} 16\% of the corpus is
title-only, half of it in \textit{Sensing}. This is the single largest recoverable loss in
the pipeline and the fix is mechanical: a scheduled retry pass over the standing
abstract-retry list, re-querying Crossref per DOI, Europe PMC and PubMed a week after
first sight, when the publisher deposit has landed. Nothing about it requires new sources.

\item \textbf{Five of 51 measurement-side records carry a health endpoint}, for the
seventh consecutive week. This remains the longest-standing structural finding in the
series, and it is now supported by 51 records in a single week rather than by a handful.

\item \textbf{Ultrafine and size-resolved number: 3 of 126.} The strongest sensing record
of the week --- a community CPC matching a research-grade reference at $R^2>0.99$ with a
3\% inter-unit coefficient of variation --- exists precisely to make that metric
affordable. The instrument is arriving faster than the studies that would use it.

\item \textbf{Two interventional records in 126}, and neither tests a monitoring
technology. The engineering turn in the indoor cluster is still being evaluated against
chambers and models rather than against health outcomes in occupied buildings.

\item \textbf{Pipeline defects carried forward}: unnormalised geography labels (which
inflate the \textit{Global / multi-region} bar in f3), the \texttt{Cancer}/\texttt{Oncologic}
split in \texttt{ENDPOINT\_CANON}, no first-author affiliation field, and a 6-element
\texttt{LIFECOURSE} array in the 2026-08-03 corpus where the schema requires 5 --- outside
this window, but it will corrupt any rollup that spans it, including the September monthly.
\end{itemize}

\clearpage

% -------------------------------------------------------------- digest
\section*{\sffamily\bfseries\color{Deep}Paper-level digest --- the week by subtopic}
\vspace{-1.5mm}{\color{Teal}\rule{\linewidth}{1.1pt}}\vspace{1mm}

{\fontsize{7.4}{9}\selectfont\color{Slate}Every record that appeared in a daily issue
keeps that issue's critical summary verbatim; nothing here is re-written or re-assessed.
Records surfaced without an abstract are shown as such. Ordered by subtopic, then by the
date of the issue that first carried the record.\par}

\input{digest_body}

% -------------------------------------------------------------- register
\section*{\sffamily\bfseries\color{Deep}Pooled corpus register}
\vspace{-1.5mm}{\color{Teal}\rule{\linewidth}{1.1pt}}\vspace{2.5mm}

{\fontsize{6.7}{8.4}\selectfont
\setlength{\tabcolsep}{2.4pt}
\begin{longtable}{@{}>{\RaggedRight}p{25mm}>{\RaggedRight}p{41mm}>{\RaggedRight}p{27mm}%
>{\RaggedRight}p{23mm}>{\RaggedRight}p{18mm}>{\RaggedRight}p{15mm}@{}}
\toprule
\textsf{\textbf{First author}} & \textsf{\textbf{Focus}} & \textsf{\textbf{Journal}} &
\textsf{\textbf{Design}} & \textsf{\textbf{Metric}} & \textsf{\textbf{Region}}\\
\midrule
\endfirsthead
\toprule
\textsf{\textbf{First author}} & \textsf{\textbf{Focus}} & \textsf{\textbf{Journal}} &
\textsf{\textbf{Design}} & \textsf{\textbf{Metric}} & \textsf{\textbf{Region}}\\
\midrule
\endhead
\multicolumn{6}{@{}l@{}}{\cellcolor{Mist}\textsf{\textbf{\textcolor{Deep}{Neuro / mental health}}}}\\[0.6mm]
\href{https://doi.org/10.1001/jamanetworkopen.2026.25015}{Dales} & Wildfire-sourced PM$_{2.5}$ and ED visits for migraine and other primary headache & \textit{JAMA Netw Open} & Case-crossover & PM$_{2.5}$ (wildfire) & Canada\\
\href{https://doi.org/10.1093/ageing/afag215}{Glover} & Umbrella meta-review of urban environmental exposures and cognitive health & \textit{Age Ageing} & Umbrella review / meta-analysis & PM$_{2.5}$ & Global\\
\href{https://doi.org/10.1016/j.nsa.2026.107019}{Goossens} & In utero air pollution exposure and pre-/postnatal brain development & \textit{Neurosci Appl} & Narrative review & PM$_{2.5}$ / PM$_{10}$ / PM$_{\mathrm{coarse}}$ & Global\\
\href{https://doi.org/10.1016/j.isci.2026.116617}{Ma} & PM$_{2.5}$-induced neurotoxicity and depression: exposure risks and mechanisms & \textit{iScience} & Narrative review & PM$_{2.5}$ + components & Global\\
\href{https://doi.org/10.3760/cma.j.cn112150-20250804-00752}{Pan} & Maternal PM exposure and offspring neurodevelopment: effects and mechanisms & \textit{Zhonghua Yu Fang Yi Xue Za Zhi} & Narrative review & PM (bulk) & China\\
\addlinespace[0.9mm]
\multicolumn{6}{@{}l@{}}{\cellcolor{Mist}\textsf{\textbf{\textcolor{Teal}{Cardiovascular \& metabolic}}}}\\[0.6mm]
\href{https://doi.org/10.3760/cma.j.cn112338-20251120-00843}{Huang} & PM$_{2.5}$ and dyslipidaemia in rural Chinese children; dietary effect modification & \textit{Zhonghua Liu Xing Bing Xue Za Zhi} & Cross-sectional (multistage) & PM$_{2.5}$ & China\\
\href{https://doi.org/10.32394/pe/218966}{Kazmierska} & Air pollution and cardiovascular disease: review of current data & \textit{Przegl Epidemiol} & Narrative review & PM$_{2.5}$ / PM$_{10}$ & Global\\
\href{https://doi.org/10.1007/s00125-026-06814-2}{Xi} & Air pollution, built environment and temperature vs. prediabetes phenotypes & \textit{Diabetologia} & Prospective cohort & PM$_{2.5}$ & Germany\\
\href{https://doi.org/10.1161/JAHA.125.046917}{Yang} & Short-term PM$_{2.5}$ and hospitalisation/ED use in atrial fibrillation & \textit{J Am Heart Assoc} & Case-crossover & PM$_{2.5}$ (wildfire-influenced) & USA\\
\addlinespace[0.9mm]
\multicolumn{6}{@{}l@{}}{\cellcolor{Mist}\textsf{\textbf{\textcolor{Sky}{Reproductive \& developmental}}}}\\[0.6mm]
\href{https://doi.org/10.1016/j.jes.2025.10.043}{Chen} & Prenatal air pollution, fetal growth trajectory phenotypes and child neurodevelopment & \textit{J Environ Sci (China)} & Prospective cohort & PM$_{2.5}$ / PM$_{10}$ & China\\
\href{https://doi.org/10.1016/j.ecoenv.2026.120566}{Huang} & Joint prenatal bisphenol + air-pollutant mixture exposure and preschool IQ & \textit{Ecotoxicol Environ Saf} & Prospective cohort (mixtures) & PM$_{2.5}$ / PM$_{10}$ & China\\
\href{https://doi.org/10.1007/s40572-026-00552-8}{Peterson} & Ambient air pollution, PAHs and pubertal development: critical review & \textit{Curr Environ Health Rep} & Critical review & PM$_{2.5}$ components & Global\\
\href{https://doi.org/10.1016/j.envres.2026.125313}{Xie} & Long-term PM$_{2.5}$ and timing of pubertal development in boys and girls & \textit{Environ Res} & Prospective cohort & PM$_{2.5}$ & China\\
\addlinespace[0.9mm]
\multicolumn{6}{@{}l@{}}{\cellcolor{Mist}\textsf{\textbf{\textcolor{Sage}{Respiratory \& allergic}}}}\\[0.6mm]
\href{https://doi.org/10.1007/s00484-026-03289-z}{Bayar Muluk} & Climate change and otolaryngology: airway disease impacts and emission reduction & \textit{Int J Biometeorol} & Narrative review & PM$_{2.5}$ & Global\\
\href{https://doi.org/10.4103/mgr.MEDGASRES-D-25-00164}{Liao} & Short-term air pollution, respiratory symptoms and lung function in asthma & \textit{Med Gas Res} & Panel study & PM$_{2.5}$ / PM$_{10}$ & China\\
\href{https://doi.org/10.3390/epidemiologia7040097}{Sanchez de Prada} & Ambient PM and asthma-coded admissions in children 0-3 y, Spain & \textit{Epidemiologia} & Ecological panel & PM$_{2.5}$ / PM$_{10}$ & Spain\\
\href{https://doi.org/10.1016/j.jhazmat.2026.143076}{Yang} & Allergenic vs. toxic risk of airborne seed hairs compared with ambient PM$_{2.5}$ & \textit{J Hazard Mater} & Field measurement & PM$_{2.5}$ + bioaerosol & China\\
\addlinespace[0.9mm]
\multicolumn{6}{@{}l@{}}{\cellcolor{Mist}\textsf{\textbf{\textcolor{Violet}{Mechanistic toxicology}}}}\\[0.6mm]
\href{https://doi.org/10.1016/j.jes.2025.12.074}{Ding} & PM$_{2.5}$ promotes vascular calcification via miR-27b-5p/HES1 signalling & \textit{J Environ Sci (China)} & In vitro + organoid + murine & PM$_{2.5}$ & China\\
\href{https://doi.org/10.1007/s00360-026-01698-5}{Maraschi} & Osmoregulatory disruption and metal bioaccumulation from settleable atmospheric PM & \textit{J Comp Physiol B} & Controlled exposure (ecotox) & Settleable APM & Brazil\\
\href{https://doi.org/10.1016/j.jes.2025.11.018}{Wu} & Macrophage NLRP3 inflammasome in real-ambient PM-induced hepatic glucose dysregulation & \textit{J Environ Sci (China)} & Real-ambient murine + co-culture & Real-ambient PM & China\\
\addlinespace[0.9mm]
\multicolumn{6}{@{}l@{}}{\cellcolor{Mist}\textsf{\textbf{\textcolor{Clay}{Occupational \& indoor}}}}\\[0.6mm]
\href{https://doi.org/10.3389/fmedt.2026.1796828}{Chen} & Nested suction instrument to cut intraoperative surgical-smoke PM$_{2.5}$ & \textit{Front Med Technol} & Bench / device experiment & PM$_{2.5}$ (surgical smoke) & China\\
\href{https://doi.org/10.1093/annweh/wxag055}{Cole} & Determinants of respirable crystalline silica exposure in Australian tunnelling & \textit{Ann Work Expo Health} & Observational exposure study & Respirable dust / RCS & Australia\\
\href{https://doi.org/10.1016/j.jhazmat.2026.143081}{Ding} & Stage-resolved hair Zn-isotope dynamics as a firefighting combustion-exposure tracer & \textit{J Hazard Mater} & Repeated-measures biomarker & Combustion-derived PM & China\\
\href{https://doi.org/10.3390/jox16040119}{Ramsden} & Xenobiotic hazards in aircraft cabin air & \textit{J Xenobiot} & Review + risk economics & Ultrafine particles & Global\\
\addlinespace[0.9mm]
\multicolumn{6}{@{}l@{}}{\cellcolor{Mist}\textsf{\textbf{\textcolor{Amber}{Sensing, forecasting \& instrumentation}}}}\\[0.6mm]
\href{https://doi.org/10.1016/j.jenvman.2026.130515}{Shi} & Adaptive spatiotemporal graph transformer for multi-site multi-step PM$_{2.5}$ forecasting & \textit{J Environ Manage} & Deep-learning model & PM$_{2.5}$ & China / India\\
\href{https://doi.org/10.3390/bios16070380}{Tsuchiya} & Seasonal and indoor-outdoor characteristics of fluorescent aerosol particles in offices & \textit{Biosensors} & Real-time sensor campaign & Bioaerosol / coarse & Japan\\
\addlinespace[0.9mm]
\multicolumn{6}{@{}l@{}}{\cellcolor{Mist}\textsf{\textbf{\textcolor{Amber}{Exposure assessment \& modelling}}}}\\[0.6mm]
\href{https://doi.org/10.1029/2025GH001621}{Solorzano-Araque} & Multiscale exposure modelling: LOTOS-EUROS + mobility surveys + particle toxicity & \textit{GeoHealth} & CTM + mobility integration & PM$_{2.5}$ / PM$_{10}$ & Colombia\\
\addlinespace[0.9mm]
\multicolumn{6}{@{}l@{}}{\cellcolor{Mist}\textsf{\textbf{\textcolor{Coral}{Burden, policy \& mitigation}}}}\\[0.6mm]
\href{https://doi.org/10.1002/hsr2.72886}{Karamzad} & Ischaemic heart disease burden in MENA 1990-2021; PM attributable fraction 28.2\% & \textit{Health Sci Rep} & GBD secondary analysis & Ambient PM & MENA\\
\href{https://doi.org/10.1371/journal.pone.0347989}{Rees} & HEPA filtration to reduce respiratory infection in care homes & \textit{PLoS One} & RCT process evaluation & Filtered PM & UK\\
\href{https://doi.org/10.1016/j.jes.2026.01.009}{Shan} & Policy-driven emission reductions and post-pandemic rebound, high-resolution inventories & \textit{J Environ Sci (China)} & Emission inventory & PM emissions & China\\
\addlinespace[0.9mm]
\multicolumn{6}{@{}l@{}}{\cellcolor{Mist}\textsf{\textbf{\textcolor{Slate}{Other clinical endpoints}}}}\\[0.6mm]
\href{https://doi.org/10.3390/dj14070440}{Lupita} & PM$_{2.5}$, preventive-care capacity and caries/edentulism burden across the EU & \textit{Dent J} & Ecological panel & PM$_{2.5}$ & EU-27\\
\href{https://doi.org/10.1016/j.jdin.2026.06.008}{Park} & Long-term PM exposure and risk of androgenetic alopecia & \textit{JAAD Int} & Nationwide cohort & PM$_{2.5}$ / PM$_{10}$ & South Korea\\
\href{https://doi.org/10.3171/2026.3.PEDS2615}{Sun} & Environmental exposure and area deprivation vs. paediatric hydrocephalus outcomes & \textit{J Neurosurg Pediatr} & Registry cohort + geospatial & PM$_{2.5}$ (satellite-derived) & USA\\
\addlinespace[0.9mm]
\bottomrule
\end{longtable}}


\vspace{4mm}
\begin{keybox}
{\sffamily\bfseries\footnotesize\color{Deep}Provenance and method}\\[1.2mm]
{\fontsize{7.2}{8.8}\selectfont
Period \textbf{2026-09-06 to 2026-09-12}, seven daily issues pooled.
\textbf{This rollup queried no literature API.} It aggregates \texttt{state/corpus/*.json},
\texttt{state/metrics.csv} and \texttt{state/trials.json} over
\texttt{PMRW\_START}--\texttt{PMRW\_END}; every record was screened, summarised and
DOI-gated at daily-issue time and is reproduced verbatim. The single exception is the
trial watch, which was queried once for the whole window against ClinicalTrials.gov with a
no-term control, because no daily window in this period had logged one.
\textbf{126 records, 126 unique DOIs} --- no duplicate publication this period. Daily
totals: 15, 20, 18, 26, 16, 10, 21. \textbf{Week-over-week comparisons} are against the
30 August--5 September rollup (150 records), with 23--29 August (109), 16--22 August (112)
and 9--15 August (134) used where a longer run is cited; the 26 July--1 August week is
excluded from "record" claims because it contains the founding retrospective back-scan.
\textbf{Effect estimates} are pooled as recorded, on heterogeneous contrasts, with no
meta-analysis and no heterogeneity statistic; several are co-pollutant rather than PM and
are labelled as such on the plot. \textbf{Figures} are regenerated over the pooled range,
not composited from the daily images. \textbf{Labels} --- subtopic, design, particle-metric,
geography --- are single-label by dominant contribution, assigned at daily-issue time, and
are a judgement rather than a controlled vocabulary.\par}
\end{keybox}

\end{document}
